\begin{question}
\noindent A space engineering student is studying satellite dish cross-sections, which are modelled using circles. Four diagrams below each illustrate a different circle theorem.

For each diagram:
\begin{itemize}
    \item[(i)] Name the circle theorem being illustrated.
    \item[(ii)] Work out the size of the missing angle marked with a letter.
\end{itemize}

\begin{center}
\begin{tabular}{cc}

\begin{tikzpicture}[scale=1.5]
    \coordinate (O) at (0,0);
    \coordinate (A) at (140:1);
    %obtuse angle B--O--A should be 116 therefore the angle between horizontal and OB must be 140-116=24
    \coordinate (B) at (24:1);
    \coordinate (C) at (250:1);
    \draw (0,0) circle (1);
    \fill[black] (O) circle (0.6pt);
    \draw (A) -- (O) -- (B);
    \draw (A) -- (C) -- (B);
    \node at (A) [above left] {$A$};
    \node at (B) [right] {$B$};
    \node at (C) [below] {$C$};
    %Going with B--O--A to ensure that the obtuse angle 116 is marked , because it is swept in the anti-clockwise
    \pic [draw, angle eccentricity=1.6, angle radius=0.5cm, "$116^\circ$"] {angle = B--O--A};
    %Going with B--C--A to ensure the correct interior angle is marked , because it is swept in the anti-clockwise
    \pic [draw, angle eccentricity=1.6, angle radius=0.4cm, "$x$"] {angle = B--C--A};
\end{tikzpicture}
&
\begin{tikzpicture}[scale=1.5]
    \coordinate (A) at (160:1);
    \coordinate (B) at (0:1);
    \coordinate (D) at (60:1);
    \coordinate (E) at (300:1);
    \draw (0,0) circle (1);
    \draw (A) -- (B);
    \draw (A) -- (D) -- (B);
    \draw (A) -- (E) -- (B);
    \node at (A) [left] {$A$};
    \node at (B) [right] {$B$};
    \node at (D) [above] {$D$};
    \node at (E) [below] {$E$};

    %use A--D--B to ensure the acute interior angle is marked, as it is swept anti-clockwise from A to B
    \pic [draw, angle eccentricity=1.5, angle radius=0.4cm, "$38^\circ$"] {angle = A--D--B};
    %use B--E--A to ensure the acute interior angle is marked, as it is swept anti-clockwise from B to A
    \pic [draw, angle eccentricity=1.5, angle radius=0.4cm, "$y$"] {angle = B--E--A};
\end{tikzpicture}
\\[2.5cm]

\begin{tikzpicture}[scale=1.5]
    \draw (0,0) circle (1);
    \fill[black] (0,0) circle (0.6pt);
    \coordinate (A) at (180:1);
    %B must be such that BOA is a straight line, therefore B should make 0 with the horizontal
    \coordinate (B) at (0:1);
    %using arbitary value for angle C
    \coordinate (C) at (125:1);
    \draw (A) -- (B);
    \draw (A) -- (C) -- (B);
    \node at (A) [left] {$A$};
    \node at (B) [right] {$B$};
    \node at (C) [above] {$C$};
    %use A--C--B to ensure the acute interior angle is marked, as it is swept anti-clockwise
    \pic [draw, angle eccentricity=1.6, angle radius=0.4cm, "$z$"] {angle = A--C--B};
\end{tikzpicture}
&
\begin{tikzpicture}[scale=1.5]
    \coordinate (A) at (140:1);
    \coordinate (B) at (60:1);
    %setting up for BAC angle 95, which means angle at center BOC = 190, the angle BOC anticlockwise would be 360-190=170. If B is 60 to the horizontal, C will be 60 + 170 = 230 to the horizontal
    \coordinate (C) at (230:1);
    \coordinate (D) at (-20:1);
    \draw (0,0) circle (1);
    \draw (A) -- (B) -- (D) -- (C) -- cycle;
    \node at (A) [left] {$A$};
    \node at (B) [above] {$B$};
    \node at (C) [below left] {$C$};
    \node at (D) [right] {$D$};
    \pic [draw, angle eccentricity=1.6, angle radius=0.4cm, "$95^\circ$"] {angle = C--A--B};
    \pic [draw, angle eccentricity=1.6, angle radius=0.4cm, "$w$"] {angle = B--D--C};
\end{tikzpicture}
\end{tabular}
\end{center}

\vspace{0.5cm}
\begin{center}
\begin{tabular}{|c|p{5.5cm}|c|}
\hline
\textbf{Diagram} & \textbf{Name of theorem} & \textbf{Missing angle} \\
\hline
1 & & $x = $ \\[1cm]
\hline
2 & & $y = $ \\[1cm]
\hline
3 & & $z = $ \\[1cm]
\hline
4 & & $w = $ \\[1cm]
\hline
\end{tabular}
\end{center}

\vspace{1cm}
\begin{flushright}
\answermarks{4}
\end{flushright}
\end{question}